\hypertarget{appendix-c80db39080be}{}
\section*{Numerical decision boundaries}
\chapter{Numerical decision boundaries}\label{appendix-c80db39080be-numerical-decision-boundaries}

\textbf{Curriculum depth:} Introduced · \textbf{First seen:} C005 · \textbf{Latest development:} C005

\section{Definition and mechanism}\label{appendix-c80db39080be-definition-and-mechanism}

A small arithmetic perturbation can cross a discrete threshold. Values just below/above pixel midpoint 2.5 choose indices 2/3 under llround; a near-zero field sign can change mesh connectivity. Scalar error and topology are different outputs.

\section{Prerequisites}\label{appendix-c80db39080be-prerequisites}

\begin{itemize}
\tightlist
\item
  \hyperref[appendix-3ae74e89b4e6]{Absolute and relative error} --- Discrepancy, magnitude and reference units define sensitivity and threshold significance.
\item
  \hyperref[appendix-637d164ae5ca]{Directed rounding modes} --- Integer/pixel rounding policy controls discrete boundary choices.
\end{itemize}

\section{Used by}\label{appendix-c80db39080be-used-by}

Connect later dependent concepts here when their lessons arrive.

\section{Recall and next depth}\label{appendix-c80db39080be-recall-and-next-depth}

Reconstruct the definition's example and name the assumptions that make it valid.

\textbf{Pending:} Robust predicates, classification uncertainty, pixel/voxel bounds and surface-extraction guarantees.

\section{Source and navigation}\label{appendix-c80db39080be-source-and-navigation}

\href{Continuous\%20Notes/Computing\%20Foundations\%20-\%20From\%20Bits\%20to\%20C++\%20and\%20Systems\%20APIs\#When\%20arithmetic\%20meets\%20uncertainty}{M005 teaching chapter} · \hyperref[appendix-fd090d982a15-m005---parts-77-88]{Part definitions} · \href{../02-Foundations-II/main.pdf\#nameddest=appendix-4ad47bdd269a}{Dependency Map (02-Foundations-II, appendix-4ad47bdd269a)}
